% ======================================================================
% CAPÍTULO 5: METODOLOGÍA EXPERIMENTAL Y BENCHMARKING (EXPANDIDO)
% ======================================================================

\chapter{Metodología Experimental y Benchmarking}
\label{cap:experimentacion}

\section{Metodología Basada en el Método Científico}
\label{sec:exp_metodo}

De acuerdo con los estándares epistemológicos de la investigación en Ciencias de la Computación y Aprendizaje Automático \cite{chalmers2013what, guillen2025libro}, la fase experimental de este Trabajo de Fin de Grado se articula en torno a las etapas formales del método científico experimental (Figura~\ref{fig:metodo_cientifico_diagrama}).

\begin{figure}[htbp]
    \centering
    \resizebox{0.92\linewidth}{!}{%
    \begin{tikzpicture}[node distance=0.7cm, auto,
        state/.style={rectangle, draw=etsiitblue, fill=blue!6, rounded corners, text width=4.2cm, text centered, minimum height=0.9cm, font=\scriptsize},
        decision/.style={diamond, draw=ugrred, fill=red!6, text width=2.0cm, text centered, inner sep=0pt, aspect=1.8, font=\scriptsize},
        conclstate/.style={rectangle, draw=darkgreen, fill=green!8, rounded corners, text width=4.2cm, text centered, minimum height=0.9cm, font=\scriptsize},
        replantstate/.style={rectangle, draw=ugrred, fill=red!6, rounded corners, text width=4.2cm, text centered, minimum height=0.9cm, font=\scriptsize},
        line/.style={draw=darkblue, -latex', thick}]
        
        % Fila 1: Planteamiento teórico inicial
        \node [state] (q) {\textbf{1. Pregunta de Investigación}\\\tiny ¿Qué arquitecturas modelan mejor bioseñales?};
        \node [state, right=0.8cm of q] (sota) {\textbf{2. Estado del Arte}\\\tiny Detección de gaps y taxonomía};
        \node [state, right=0.8cm of sota] (hip) {\textbf{3. Hipótesis Científicas}\\\tiny Formulación formal $H_1, H_2, H_3$};
        
        % Fila 2: Fase experimental y análisis
        \node [state, below=0.8cm of hip] (exp) {\textbf{4. Diseño Experimental}\\\tiny GroupKFold + Optuna + Clúster GPU};
        \node [state, below=0.8cm of sota] (res) {\textbf{5. Análisis de Resultados}\\\tiny Métricas $R^2$, MAE, RMSE, Pearson $r$};
        
        % Fila 3: Evaluación crítica y bifurcación
        \node [decision, below=0.8cm of res] (val) {¿Hipótesis\\confirmadas?};
        \node [conclstate, left=0.8cm of val] (concl) {\textbf{6. Conclusiones y ODS}\\\tiny Transferencia e impacto};
        \node [replantstate, right=0.8cm of val] (replantear) {\textbf{Reevaluación Crítica}\\\tiny Estudio de limitaciones};
        
        % Conexiones del flujo principal
        \path [line] (q) -- (sota);
        \path [line] (sota) -- (hip);
        \path [line] (hip) -- (exp);
        \path [line] (exp) -- (res);
        \path [line] (res) -- (val);
        
        % Bifurcación
        \path [line] (val) -- node [above, font=\scriptsize\bfseries] {Sí} (concl);
        \path [line] (val) -- node [above, font=\scriptsize\bfseries] {No} (replantear);
        
        % Bucle de retroalimentación ortogonal sin cruces
        \path [line] (replantear.east) -- ++(0.5,0) |- node [near end, right, font=\tiny\bfseries] {Reformulación} (hip.east);
    \end{tikzpicture}%
    }
    \caption{Ciclo metodológico del método científico aplicado a la experimentación del TFG: flujo secuencial riguroso y retroalimentación ortogonal.}
    \label{fig:metodo_cientifico_diagrama}
\end{figure}

\subsection{Formulación de Hipótesis Científicas}

A partir de la revisión del estado del arte realizada en el Capítulo~\ref{cap:estado_del_arte}, se plantean tres hipótesis de trabajo formales:

\begin{hipotesis}[\textbf{Superioridad del Modelado Temporal Profundo}]
\label{hip:h1}
Las redes neuronales profundas que modelan explícitamente la dinámica temporal secuencial (LSTM, GRU, TCN, \xlstm{}) presentan una capacidad de generalización sobre participantes no vistos significativamente superior a la de los modelos clásicos de Machine Learning basados en características tabulares agregadas.
\end{hipotesis}

\begin{hipotesis}[\textbf{Estabilidad y Eficacia de la Celda sLSTM de xLSTM}]
\label{hip:h2}
El mecanismo de compuertas exponenciales y normalización logarítmica ($m_t, n_t$) de la celda estabilizada \slstm{} permite una convergencia más rápida y un menor error cuadrático medio en la regresión de fatiga multicanal frente a la arquitectura clásica LSTM con compuertas sigmoides.
\end{hipotesis}

\begin{hipotesis}[\textbf{Capacidad y Limitaciones de los Modelos Fundacionales}]
\label{hip:h3}
Los Modelos Fundacionales preentrenados a gran escala (\chronos{}, \timesfm{}, \moment{}) demuestran capacidades predictivas viables en régimen \textit{zero-shot} sobre bioseñales, pero su precisión bruta es superada por modelos recurrentes ligeros optimizados bayesianamente sobre el dominio específico, incurriendo además los primeros en un coste computacional sensiblemente mayor.
\end{hipotesis}


% ======================================================================
\section{Protocolo de Validación Cruzada GroupKFold por Participante}
\label{sec:exp_groupkfold}
% ======================================================================

Uno de los errores metodológicos más graves y recurrentes en la literatura de bioseñales y computación vestible es la \textbf{fuga de datos inter-sujeto (\textit{subject data leakage})} \cite{kalanadhabhatta2022fatigueset}. 

Si las ventanas temporales generadas a partir de las sesiones de un mismo sujeto se distribuyen aleatoriamente entre el conjunto de entrenamiento y el conjunto de validación (mediante un \textit{K-Fold} convencional o división aleatoria simple), el modelo neuronal tiende a memorizar los patrones fisiológicos basales, la morfología específica del ECG o la impedancia de la piel del individuo en lugar de aprender representaciones universales de la fatiga. Esto conduce a un sobreajuste severo y a métricas de rendimiento ilusoriamente optimistas que colapsan al aplicar el sistema sobre nuevos usuarios en la práctica real.

Para erradicar de forma estricta este problema, en este trabajo se establece un protocolo inmutable de \textbf{Validación Cruzada Estratificada por Grupos (\textit{GroupKFold Cross-Validation})} con $K=5$ particiones:
\begin{equation}
    \mathcal{D} = \bigcup_{k=1}^5 \mathcal{F}_k, \quad \mathcal{S}(\mathcal{F}_i) \cap \mathcal{S}(\mathcal{F}_j) = \emptyset \quad \forall i \ne j
\end{equation}
donde $\mathcal{S}(\mathcal{F}_k)$ representa el conjunto de identificadores de participantes asignados al fold $k$. En cada iteración $k$, los modelos se entrenan exclusivamente sobre las ventanas de los participantes en $\mathcal{D} \setminus \mathcal{F}_k$ y se evalúan a ciegas sobre los participantes inéditos de $\mathcal{F}_k$. La Tabla~\ref{tab:groupkfold_splits} detalla la partición exacta de participantes y la distribución porcentual de muestras entre conjuntos en cada uno de los 5 pliegues de evaluación.

\begin{table}[htbp]
\centering
\caption{Asignación de participantes y distribución de particiones en GroupKFold ($K=5$).}
\label{tab:groupkfold_splits}
\small
\begin{tabularx}{\linewidth}{@{}ccXX@{}}
\toprule
\textbf{Fold} & \textbf{Participantes en Test/Val} & \textbf{Participantes en Train} & \textbf{Distribución Muestral} \\
\midrule
Fold 1 & P01, P02 & P03, P04, P05, P06, P07, P08, P09, P10, P11, P12 & 18.2\% Val / 81.8\% Train \\
Fold 2 & P03, P04 & P01, P02, P05, P06, P07, P08, P09, P10, P11, P12 & 17.5\% Val / 82.5\% Train \\
Fold 3 & P05, P06 & P01, P02, P03, P04, P07, P08, P09, P10, P11, P12 & 19.1\% Val / 80.9\% Train \\
Fold 4 & P07, P08 & P01, P02, P03, P04, P05, P06, P09, P10, P11, P12 & 21.0\% Val / 79.0\% Train \\
Fold 5 & P09, P10, P11, P12 & P01, P02, P03, P04, P05, P06, P07, P08 & 24.2\% Val / 75.8\% Train \\
\bottomrule
\end{tabularx}
\end{table}


% ======================================================================
\section{Métricas de Evaluación Multiobjetivo}
\label{sec:exp_metricas}
% ======================================================================

Dado que el problema se formula como una regresión multivariada continua de dos dimensiones $\bm{y} = [y_{\text{física}}, y_{\text{mental}}]^T \in [0, 100]^2$, se calculan las siguientes métricas para cada dimensión $d \in \{\text{física}, \text{mental}\}$ sobre los $N$ ejemplos del conjunto de validación:

\subsection{Error Absoluto Medio (MAE)}
Cuantifica la magnitud media de los errores en la escala original de la variable (puntos VAS de 0 a 100):
\begin{equation}
    \text{MAE}^{(d)} = \frac{1}{N} \sum_{i=1}^N \left| y_i^{(d)} - \hat{y}_i^{(d)} \right|
\end{equation}

\subsection{Raíz del Error Cuadrático Medio (RMSE)}
Penaliza de forma cuadrática las desviaciones de gran magnitud, siendo un indicador clave de robustez ante fallos atípicos:
\begin{equation}
    \text{RMSE}^{(d)} = \sqrt{\frac{1}{N} \sum_{i=1}^N \left( y_i^{(d)} - \hat{y}_i^{(d)} \right)^2}
\end{equation}

\subsection{Coeficiente de Determinación ($R^2$)}
Evalúa la proporción de la varianza de la variable dependiente explicada por el modelo predictivo frente al estimador trivial de la media $\bar{y}^{(d)}$:
\begin{equation}
    R^{2\,(d)} = 1 - \frac{\sum_{i=1}^N \left( y_i^{(d)} - \hat{y}_i^{(d)} \right)^2}{\sum_{i=1}^N \left( y_i^{(d)} - \bar{y}^{(d)} \right)^2}
\end{equation}
En validación cruzada inter-sujeto con alta variabilidad, valores de $R^2 < 0$ indican que el error cuadrático del modelo es superior a la varianza de los datos en ese fold específico, un fenómeno ampliamente documentado en problemas biomédicos complejos con cohortes reducidas \cite{kalanadhabhatta2022fatigueset}.

\subsection{Coeficiente de Correlación de Pearson ($r$)}
Mide la correlación lineal entre los valores reales y las predicciones, evaluando si el modelo captura la tendencia ascendente o descendente de la fatiga con independencia de sesgos de escala global:
\begin{equation}
    r^{(d)} = \frac{\sum_{i=1}^N (y_i^{(d)} - \bar{y}^{(d)})(\hat{y}_i^{(d)} - \bar{\hat{y}}^{(d)})}{\sqrt{\sum_{i=1}^N (y_i^{(d)} - \bar{y}^{(d)})^2} \sqrt{\sum_{i=1}^N (\hat{y}_i^{(d)} - \bar{\hat{y}}^{(d)})^2}}
\end{equation}

\subsection{Métricas Probabilísticas (para Modelos Fundacionales)}
Para los modelos que generan distribuciones predictivas completas (como \chronos{}), se evalúa el \textbf{Continuous Ranked Probability Score (CRPS)}:
\begin{equation}
    \text{CRPS}(F, y) = \int_{-\infty}^\infty \left( F(z) - \mathbb{I}(z \ge y) \right)^2 dz
\end{equation}
junto con la cobertura empírica del intervalo de predicción al 90\% ($PICP_{90}$):
\begin{equation}
    PICP_{90} = \frac{1}{N}\sum_{i=1}^N \mathbb{I}\left( \hat{y}_{i}^{0.05} \le y_i \le \hat{y}_{i}^{0.95} \right) \times 100\%
\end{equation}


% ======================================================================
\section{Infraestructura de Computación en Servidor GPU y SLURM}
\label{sec:exp_cluster}
% ======================================================================

El entrenamiento y benchmarking a gran escala de las arquitecturas profundas y los Modelos Fundacionales (cuyos tamaños oscilan entre 200 y 710 millones de parámetros) exige recursos computacionales de alto rendimiento.

\subsection{Especificaciones del Clúster GPU de la Universidad de Granada}
Los experimentos se ejecutaron en el nodo de computación \code{titan} del clúster de altas prestaciones de la UGR:
\begin{itemize}
    \item \textbf{Nodo de Acceso (Cabecera):} \code{ngpu.ugr.es} bajo GNU/Linux x86\_64.
    \item \textbf{Aceleración Gráfica:} GPUs dedicadas NVIDIA con 12 GB de memoria VRAM de alta velocidad.
    \item \textbf{Memoria de Sistema y Almacenamiento:} 32 GB de RAM asignados por \textit{job} en partición \code{dios} y almacenamiento persistente de 239 TB en red en \code{/mnt/homeGPU/\-egullonl01/}.
    \item \textbf{Gestor de Carga de Trabajo:} SLURM (\textit{Simple Linux Utility for Resource Management}).
\end{itemize}

\subsection{Resolución de Desafíos de Ingeniería en el Clúster}
Durante la fase de despliegue se identificaron y subsanaron tres retos operacionales críticos:
\begin{enumerate}
    \item \textbf{Superación de la Cuota de Disco en Directorio de Usuario:} Por defecto, las librerías HuggingFace y los runtimes de Jupyter escriben en \code{\textasciitilde{}/.cache/}, provocando el error \textit{Disk quota exceeded} debido al límite estricto de 5 GB en el directorio \code{HOME}. Se solucionó redirigiendo las variables de entorno de caché al almacenamiento persistente de red:
\begin{lstlisting}[language=bash,breaklines=true,basicstyle=\footnotesize\ttfamily]
export HF_HOME="/mnt/homeGPU/egullonl01/.cache/huggingface"
export XDG_CACHE_HOME="/mnt/homeGPU/egullonl01/.cache"
export IPYTHONDIR="/mnt/homeGPU/egullonl01/.ipython"
export JUPYTER_RUNTIME_DIR=\
  "/mnt/homeGPU/egullonl01/.local/share/jupyter/runtime"
\end{lstlisting}
    \item \textbf{Fugas de Memoria VRAM entre Modelos:} PyTorch no devuelve inmediatamente la memoria reservada al driver mientras el proceso principal siga activo. Para evitar errores \textit{CUDA Out of Memory} (OOM) en benchmarks continuos, se programó la ejecución secuencial de cada familia de modelos en procesos limpios e independientes mediante \code{jupyter nbconvert --execute --inplace}.
    \item \textbf{Límites de Tiempo en SLURM (12h máximo):} Se implementó un sistema de persistencia en disco de puntos de control (\textit{checkpoints}) por época y almacenamiento en base de datos SQLite (\code{optuna\_v2.db}) con soporte multi-hilo para permitir la reanudación transparente de estudios interrumpidos.
\end{enumerate}


% ======================================================================
\section{Diseño de las Cuatro Fases Experimentales}
\label{sec:exp_fases}
% ======================================================================

El plan experimental se organizó en cuatro fases progresivas de complejidad creciente:
\begin{itemize}
    \item \textbf{Fase 1 (Baselines Clásicos de Machine Learning):} Evaluación de Random Forest (100 árboles), K-Nearest Neighbors ($k=3$), Árbol de Decisión, Ridge y SVR sobre vectores de características agregadas por ventana.
    \item \textbf{Fase 2 (Redes Neuronales Recurrentes Base):} Evaluación de RNN simple, LSTM clásica y GRU con optimización de hiperparámetros.
    \item \textbf{Fase 3 (Arquitecturas Avanzadas de Series Temporales):} Búsqueda bayesiana con Optuna sobre CNN-LSTM, TCN (convoluciones causales dilatadas), Transformer, \patchtst{} y la nueva arquitectura \xlstm{} (celda \slstm{} estabilizada).
    \item \textbf{Fase 4 (Modelos Fundacionales de Series Temporales):} Evaluación comparativa de \moment{}-1-large (fine-tuning de cabeza lineal), \chronos{}-t5-base (zero-shot probabilístico) y \timesfm{} 2.5 (zero-shot causal con probe de salida).
\end{itemize}
